\documentclass[a4paper]{article}

\usepackage{graphicx}
\usepackage{amsmath,amssymb}
\usepackage{hyperref}
\usepackage{minted}
\usepackage{multirow}
\usepackage{comment}
\usepackage{float}
\usepackage{todonotes}

\title{Experimental setup}
\author{}
\date{}

\begin{document}
   \maketitle
   This part provides a template for describing the conditions under which the experiment was carried out.




   \section{Goal}
      The goal is for two robots to learn how to move around a square-like corridor without colliding with the walls under human teleportation control, and then to compare this experience with experiences in which they must perform maneuvers to avoid obstacles caused by changes made to the corridor.
   \section{Involved robots}
      During online telemetry recording, a leader--follower strategy is used by the two UAVs involved in the environment. The leader is the robot to which the human operator sends control inputs. The follower follows a set of rules that depend on the leader's positional state. As described later, the human and the environment are also considered robots.
      \subsection{Simulation environment}
         \input{/nd_research_publication/src/kind/journal_2026/experiments/simulation/simulation.tex}
      \subsection{UAV prototype}
         Tarot T650







         \subsubsection{Sensors}
            % to find frequencies, run /home/donkarlo/Dropbox/repo/nd_robotic_project/src/nd_robotic_ai/platform/ros/message/message_frequency_finder.py

            \paragraph{GPS}
               \begin{itemize}
                  \item \textbf{Frequency:} 83.33 Hz
               \end{itemize}

            \paragraph{LIDAR}
               \begin{itemize}
                  \item \textbf{Frequency:} 20.83 Hz
                  \item 720 dimensions.
               \end{itemize}
               In each scan, the 2D RP LiDAR A2 produces 720 distance measurements, each separated from the next by 0.14 degrees, and scans as far as 14 km. This produces 720 points. In this research, every distance greater than 4 m is treated as 4 m so that the wall on the opposite side of the corridor is omitted from the UAV's perception. DBSCAN clustering is performed on the polygon formed by connecting each pair of consecutive points.
               
               

      \subsection{Environment}
         The world is simulated in Gazebo, and the robot simulation is implemented using ROS.

      \subsection{Human}
         The human is also considered a robot.

      \subsection{Homogeneity}
         Homogeneous

      \subsection{Dependency}
         Leader--follower
   \section{Scenarios}
   \include{/nd_research_publication/src/kind/journal_2026/experiments/comparative/condition/normal/normal.tex}
   \include{/home/donkarlo/Dropbox/repo/nd_research_publication/src/kind/journal_2026/experiments/comparative/condition/follow/follow.tex}

\end{document}